\chapter*{Conclusion Générale et Perspectives}
\addcontentsline{toc}{chapter}{Conclusion Générale et Perspectives} % inclusion dans la table des matières

Le présent rapport a étudié la \textbf{détection de fraude bancaire avant et après l’apparition du Big Data}. 
Nous avons montré que l’évolution technologique et la disponibilité de grandes quantités de données ont profondément transformé les méthodes de détection, 
en permettant une analyse plus rapide, précise et évolutive.

Les principaux points retenus sont :  
\begin{enumerate}
    \item Les méthodes traditionnelles basées sur des règles statiques restent limitées face aux fraudes complexes et aux volumes croissants de transactions.  
    \item L’émergence du Big Data a permis de développer des systèmes plus performants, exploitant des algorithmes d’apprentissage automatique et des analyses en temps réel.  
    \item La conception du système de détection s’appuie sur des diagrammes UML pour modéliser les acteurs, les processus et les interactions, assurant ainsi une structure claire et évolutive.  
    \item La visualisation des résultats via des tableaux de bord interactifs facilite la compréhension et le suivi des alertes par les analystes.  
\end{enumerate}

Ce travail a permis d’allier les connaissances théoriques acquises en cours avec des applications pratiques, 
en offrant une approche complète du problème de la fraude bancaire et de ses solutions modernes.
\clearpage
\section{Perspectives}
Pour améliorer et prolonger ce travail, plusieurs pistes peuvent être envisagées :  
\begin{enumerate}
    \item \textbf{Intégration de données en temps réel} : connecter le système directement aux flux transactionnels pour détecter la fraude instantanément.  
    \item \textbf{Amélioration des modèles de détection} : tester et comparer différents algorithmes d’apprentissage automatique et de Deep Learning pour augmenter la précision.  
    \item \textbf{Extension des visualisations} : créer des dashboards plus interactifs, avec des filtres dynamiques et des alertes automatiques.  
    \item \textbf{Analyse prédictive} : utiliser les données historiques pour anticiper les fraudes futures et établir des stratégies proactives.  
    \item \textbf{Sécurité et confidentialité} : renforcer la protection des données sensibles et appliquer les normes réglementaires (RGPD, etc.).  
\end{enumerate}

Ces perspectives montrent que la combinaison du Big Data, de l’intelligence artificielle et de la visualisation 
constitue une voie prometteuse pour améliorer la sécurité et l’efficacité dans le secteur bancaire.
